\subsection{观察者索引、信息状态与 forcing 语义}\label{subsec:observer-indexed-forcing}
前述"地址先于取值""\(\mathsf{NULL}\) 的局部截面判据"与"多轴前缀提升可读性"本质上都属于对象层可断言性的组织问题。
为把这三条链压缩到同一语义框架中，本小节把对象层事实统一放回\emph{观察者索引的信息状态}上：全局层只保留推理规则，而对象层命题是否成立，则由局部信息状态上的 forcing 关系给出。

\begin{definition}[观察者索引的信息状态框架]\label{def:observer-indexed-information-frame}
称八元组
\[
\mathfrak F=(P,\le,\iota,\tau,\mathrm{Cov},\mathrm{Upd},D,V)
\]
为一个观察者索引的信息状态框架，若满足下列条件：
\begin{enumerate}
  \item \(P\) 为信息状态集合；每个 \(p\in P\) 表示一个局部可访问的信息条件，而非无索引的全局世界。
  \item \(\le\) 为 \(P\) 上的预序；若 \(q\le p\)，则称 \(q\) 是 \(p\) 的细化，表示 \(q\) 在 \(p\) 的基础上拥有更多记录、更细的区分能力或更强的分类结构。
  \item \(\iota:P\to I\) 为观察者索引映射，\(\tau:P\to T\) 为时间索引映射，并满足
  \[
  q\le p \Longrightarrow \iota(q)=\iota(p),\qquad \tau(q)\ge \tau(p).
  \]
  \item 对每个 \(p\in P\)，\(\mathrm{Cov}(p)\) 给出 \(p\) 的覆盖族；每个覆盖
  \[
  \{p_k\to p\}_{k\in K}
  \]
  满足 \(p_k\le p\)，并且覆盖在细化下稳定：若 \(q\le p\) 且 \(\{p_k\to p\}_{k\in K}\in\mathrm{Cov}(p)\)，则存在 \(\{q_\ell\to q\}_{\ell\in L}\in\mathrm{Cov}(q)\) 及映射 \(\lambda:L\to K\)，使得 \(q_\ell\le p_{\lambda(\ell)}\) 对一切 \(\ell\in L\) 成立。
  \item 对每个更新动作 \(\alpha\)，\(\mathrm{Upd}_\alpha(p)\subseteq P\) 为在 \(p\) 上执行 \(\alpha\) 后可到达的合法后继状态，并满足与细化的双向相容性：若 \(q\le p\)，则
  \[
  r\in \mathrm{Upd}_\alpha(q)\Rightarrow \exists s\in \mathrm{Upd}_\alpha(p)\ \text{使}\ r\le s,
  \]
  \[
  s\in \mathrm{Upd}_\alpha(p)\Rightarrow \exists r\in \mathrm{Upd}_\alpha(q)\ \text{使}\ r\le s.
  \]
  \item \(D\) 为对象域；对每个原子命题 \(A\)，赋值 \(V(A)\subseteq P\) 为下闭集，即
  \[
  p\in V(A),\ q\le p \Longrightarrow q\in V(A).
  \]
\end{enumerate}
\end{definition}

\begin{definition}[forcing 关系]\label{def:observer-indexed-forcing}
在框架 \(\mathfrak F\) 上，定义对象层可断言关系
\[
p\Vdash \varphi
\]
如下：
\begin{enumerate}
  \item 对原子命题 \(A\)，规定
  \[
  p\Vdash A \Longleftrightarrow p\in V(A).
  \]
  \item 对合取，规定
  \[
  p\Vdash \varphi\land\psi
  \Longleftrightarrow
  p\Vdash \varphi\ \text{且}\ p\Vdash \psi.
  \]
  \item 对蕴含，规定
  \[
  p\Vdash \varphi\to\psi
  \Longleftrightarrow
  \forall q\le p,\ \bigl(q\Vdash \varphi \Rightarrow q\Vdash \psi\bigr).
  \]
  \item 对否定，规定
  \[
  p\Vdash \neg\varphi
  \Longleftrightarrow
  \forall q\le p,\ q\nVdash \varphi.
  \]
  \item 对析取，规定 \(p\Vdash \varphi\lor\psi\) 当且仅当存在某个覆盖 \(\{p_k\to p\}_{k\in K}\in\mathrm{Cov}(p)\)，使得对每个 \(k\in K\)，均有
  \[
  p_k\Vdash \varphi\quad \text{或}\quad p_k\Vdash \psi.
  \]
  \item 对全称量词，规定
  \[
  p\Vdash \forall x\,\varphi(x)
  \Longleftrightarrow
  \forall q\le p,\ \forall d\in D,\ q\Vdash \varphi(d).
  \]
  \item 对存在量词，规定 \(p\Vdash \exists x\,\varphi(x)\) 当且仅当存在某个覆盖 \(\{p_k\to p\}_{k\in K}\in\mathrm{Cov}(p)\)，并可为每个 \(k\in K\) 选取 \(d_k\in D\)，使得
  \[
  p_k\Vdash \varphi(d_k).
  \]
  \item 对更新模态，规定
  \[
  p\Vdash [\alpha]\varphi
  \Longleftrightarrow
  \forall q\in \mathrm{Upd}_\alpha(p),\ q\Vdash \varphi,
  \]
  \[
  p\Vdash \langle\alpha\rangle\varphi
  \Longleftrightarrow
  \exists q\in \mathrm{Upd}_\alpha(p),\ q\Vdash \varphi.
  \]
\end{enumerate}
\end{definition}

\begin{definition}[已决定]\label{def:observer-indexed-decided}
称命题 \(\varphi\) 在状态 \(p\) 上\emph{已决定}，若
\[
p\Downarrow \varphi
\quad:\Longleftrightarrow\quad
p\Vdash \varphi\ \text{或}\ p\Vdash \neg\varphi.
\]
若既无 \(p\Vdash \varphi\)，又无 \(p\Vdash \neg\varphi\)，则称 \(\varphi\) 在 \(p\) 上未决定。
\end{definition}

\begin{proposition}[forcing 的单调性]\label{prop:observer-indexed-forcing-monotonicity}
若 \(q\le p\) 且 \(p\Vdash \varphi\)，则 \(q\Vdash \varphi\)。
\end{proposition}
\leanverified{paper\_recursive\_addressing\_observer\_indexed\_forcing\_monotonicity}

\begin{proof}
对公式复杂度作归纳。
原子情形由 \(V(A)\) 的下闭性直接得到；合取情形平凡。
蕴含、否定与全称量词情形都只需注意：定义中对一切 \(r\le p\) 的量化在 \(q\le p\) 下仍有效。
析取与存在量词情形由覆盖在细化下的稳定性推出：\(p\) 上的任一合法覆盖都可沿 \(q\le p\) 限制为 \(q\) 的合法覆盖。
对更新模态，\([\alpha]\) 的情形由后继状态与细化的前向相容性推出；\(\langle\alpha\rangle\) 的情形则由后继状态与细化的后向相容性推出。
故结论成立。
\end{proof}

\begin{corollary}[细化不推翻已建立的对象事实]\label{cor:observer-indexed-no-retraction}
若某命题在较粗状态 \(p\) 上已可断言，则在任何后续细化状态 \(q\le p\) 上仍可断言。
因此，信息细化的作用不是撤销对象层事实，而是缩小未决定域。
\end{corollary}
\leanverified{paper\_recursive\_addressing\_observer\_indexed\_no\_retraction}

\begin{proof}
这是命题 \ref{prop:observer-indexed-forcing-monotonicity} 的直接推论。
\end{proof}

\begin{proposition}[全局层只统一推理规则]\label{prop:observer-indexed-global-layer-logic-only}
设 \(\Gamma\vdash \varphi\) 是与定义 \ref{def:observer-indexed-forcing} 相容的演算层推理。
则对任意状态 \(p\in P\)，若
\[
p\Vdash \gamma
\qquad (\forall \gamma\in \Gamma),
\]
则
\[
p\Vdash \varphi.
\]
因此，全局层统一的是推理规则，而对象层事实始终由具体信息状态上的 forcing 给出。
\end{proposition}
\leanverified{paper\_recursive\_addressing\_observer\_indexed\_global\_layer\_logic\_only}

\begin{proof}
这是 forcing 语义对相应演算层的标准可靠性陈述；证明对推导长度作归纳即可。
证毕。
\end{proof}

\begin{definition}[地址对象与局部截面]\label{def:observer-indexed-address-local-sections}
沿用定义 \ref{def:prefix-site}--\ref{def:prefix-sheaf-readout} 的前缀站点与协议可接受读出层。
对每个信息状态 \(p\)，记
\[
A(p)\subseteq \mathcal{A}_{\mathrm{adm}}
\]
为在 \(p\) 上已对象化且协议允许的地址集合。
对每个 \(a\in A(p)\)，记
\[
\Gamma_p(a)\subseteq \mathscr{F}(a)
\]
为在状态 \(p\) 上、地址 \(a\) 处全部协议允许、类型意义下单值且可核验的局部截面之集合。
若 \(a\notin A(p)\)，则称该地址在 \(p\) 上尚未对象化；若 \(a\in A(p)\) 但 \(\Gamma_p(a)=\varnothing\)，则称该地址在 \(p\) 上尚无合法局部读出。
\end{definition}

\begin{proposition}[地址先于取值]\label{prop:observer-indexed-address-before-value}
对任意状态 \(p\) 与地址符号 \(a\)，若 \(a\notin A(p)\)，则一切形如"\(a\) 的值为 \(v\)"的对象层命题在 \(p\) 上都不能成立；
若 \(a\in A(p)\) 但 \(\Gamma_p(a)=\varnothing\)，则"\(a\) 在当前类型中可被单值读出"的命题在 \(p\) 上不能成立。
\end{proposition}
\leanverified{paper\_observer\_indexed\_address\_before\_value\_seeds}
\leanverified{paper\_observer\_indexed\_address\_before\_value\_package}
\leanverified{paper\_observer\_indexed\_address\_before\_value}

\begin{proof}
对象层取值命题要成立，前提是存在合法指称对象与相应的局部截面。
若地址对象尚未形成，则命题缺乏对象语义支撑；若地址已对象化但局部截面为空，则读出命题缺乏证书支撑。
因此两种情形下相关命题都不能在 \(p\) 上被强制。
\end{proof}

\begin{corollary}[非 \texorpdfstring{$\mathsf{NULL}$}{NULL} 的 forcing 判据]\label{cor:observer-indexed-nonnull-criterion}
在状态 \(p\) 上，地址 \(a\) 的读出为非 \(\mathsf{NULL}\) 当且仅当
\[
a\in A(p)
\qquad\text{且}\qquad
\Gamma_p(a)\neq \varnothing.
\]
因此，非 \(\mathsf{NULL}\) 不是一个值域元素，而是局部截面的存在性条件。
\end{corollary}
\leanverified{paper\_recursive\_addressing\_observer\_indexed\_nonnull\_criterion}

\begin{proof}
由命题 \ref{prop:observer-indexed-address-before-value} 立即得到。
\end{proof}

\begin{definition}[\texorpdfstring{$\mathsf{NULL}$}{NULL} 的状态化三分解]\label{def:observer-indexed-null-trichotomy}
在状态 \(p\) 上，地址 \(a\) 的 \(\mathsf{NULL}\) 分为三类：
\begin{enumerate}
  \item \textbf{语义空值：} \(a\notin A(p)\)。此时对象层尚未形成合法指称。
  \item \textbf{协议空值：} \(a\in A(p)\) 但由于当前可见域、类型约束或协议纪律的限制，\(\Gamma_p(a)=\varnothing\)。此时有地址而无合法读出。
  \item \textbf{碰撞空值：} 地址与协议均合法，但底层读出仍存在不可消解的多对一压缩；在未补足足够残差轴、未提升到单射层级之前，\(\Gamma_p(a)=\varnothing\)。
\end{enumerate}
\end{definition}

\begin{proposition}[\texorpdfstring{$\mathsf{NULL}$}{NULL} 的结构性]\label{prop:observer-indexed-null-structural}
\(\mathsf{NULL}\) 不是值域中的一个普通元素，而是对象层 forcing 失败的结构性形式。
\leanverified{paper\_recursive\_addressing\_observer\_indexed\_null\_structural}
\end{proposition}

\begin{proof}
由推论 \ref{cor:observer-indexed-nonnull-criterion}，非 \(\mathsf{NULL}\) 的充要条件是局部截面的非空存在。
故 \(\mathsf{NULL}\) 的本质不是某个特殊值被读出，而是相关读出命题未能在当前状态上被强制成立。
\end{proof}

\begin{definition}[多轴前缀与公共细化]\label{def:observer-indexed-common-refinement}
设 \(u\) 与 \(v\) 是两个协议兼容的前缀轴，并记 \(p_u,p_v\in P\) 为分别只包含单轴 \(u,v\) 信息的状态。
若存在状态 \(p_{u\otimes v}\in P\) 使
\[
p_{u\otimes v}\le p_u,
\qquad
p_{u\otimes v}\le p_v,
\]
则称 \(p_{u\otimes v}\) 为它们的联合前缀状态或公共细化。
\end{definition}

\begin{proposition}[多轴前缀提升可断言性]\label{prop:observer-indexed-common-refinement-decidability}
\leanverified{paper\_recursive\_addressing\_observer\_indexed\_common\_refinement\_decidability}
在定义 \ref{def:observer-indexed-common-refinement} 的前提下，联合前缀不会降低任何已建立的对象层事实。
并且，凡依赖于额外残差轴才能消解的碰撞命题，都可能在公共细化状态上首次变为已决定。
\end{proposition}

\begin{proof}
第一部分是命题 \ref{prop:observer-indexed-forcing-monotonicity} 的直接结果。
第二部分则来自公共细化对不可区分候选读出的进一步切分：原先在单轴状态上混合在一起的候选项，经附加轴区分后，可能首次形成单值且可核验的局部截面，从而把未决定命题推进到已决定状态。
\end{proof}

\begin{definition}[延迟分类]\label{def:observer-indexed-delayed-classification}
设 \(s\le \tau(p)\)，且 \(A_s\) 为一个关于较早时刻 \(s\) 的命题。
若在状态 \(p\) 上
\[
p\nVdash A_s,
\qquad
p\nVdash \neg A_s,
\]
但存在某个更新动作 \(\alpha\) 与某个后继状态 \(q\in \mathrm{Upd}_\alpha(p)\)，使得
\[
q\Downarrow A_s,
\]
则称 \(A_s\) 在 \(p\to q\) 的更新过程中发生了延迟分类。
\end{definition}

\begin{proposition}[延迟分类不改写过去]\label{prop:observer-indexed-delayed-classification-not-retrocausal}
在定义 \ref{def:observer-indexed-delayed-classification} 的情形下，更新所改变的是关于时刻 \(s\) 之命题的当前可断言性，而不是时刻 \(s\) 的对象层内容本身。
\end{proposition}

\begin{proof}
命题 \(A_s\) 的时间索引 \(s\) 在更新前后保持不变；变化的只是信息状态由 \(p\) 细化为某个可达后继 \(q\)。
因此，延迟分类表达的是"关于过去对象的命题在更强信息状态上首次被定解"，而非"过去被回写为另一个对象层事实"。
\end{proof}
\leanverified{paper\_recursive\_addressing\_observer\_indexed\_delayed\_classification\_not\_retrocausal}

\begin{corollary}[观察更新的对象层意义]\label{cor:observer-indexed-update-decides-not-create}
观察、再分类、增加残差轴、拼接证书对象等操作，在对象层上的统一作用，是把某些命题从未决定推进到已决定，而不是凭空创造事实。
\leanverified{paper\_recursive\_addressing\_observer\_indexed\_update\_decides\_not\_create}
\end{corollary}

\begin{proposition}[递归柱事件的有限构造证书与 incidence 账本]\label{prop:distill-euclid-elements-recursive-cylinder-construction-certificate}
\textbf{Dependency status: construction.}
沿用命题 \ref{prop:recursive-addressing-refinement} 的记号。给定
\[
a=(a_0,\ldots,a_{\ell-1})\in\bigl(\{0,1\}^{r_L}\bigr)^\ell,
\qquad
a_j=(a_{j,1},\ldots,a_{j,r_L}),
\]
定义有符号原子
\[
E_{j,i}(a):=
\begin{cases}
T^{-(t+j)}C_i^{(L)},& a_{j,i}=1,\\
X\setminus T^{-(t+j)}C_i^{(L)},& a_{j,i}=0.
\end{cases}
\]
则有限列表
\[
\mathfrak C(L,a,t):=\bigl(E_{j,i}(a)\bigr)_{0\le j<\ell,\ 1\le i\le r_L}
\]
是 \(C^{(L+1)}_{a,t}\) 的构造证书，并且其 incidence 账本满足：
\begin{enumerate}
  \item 精确构造式
  \[
  C^{(L+1)}_{a,t}
  =\bigcap_{j=0}^{\ell-1}\bigcap_{i=1}^{r_L}E_{j,i}(a);
  \]
  \item 若 \(b=(b_0,\ldots,b_{\ell'-1})\) 满足 \(\ell'\ge\ell\) 且 \(b_j=a_j\) 对 \(0\le j<\ell\) 成立，则
  \[
  C^{(L+1)}_{b,t}\subseteq C^{(L+1)}_{a,t};
  \]
  \item 若 \(a,a'\in(\{0,1\}^{r_L})^\ell\) 且存在 \((j,i)\) 使 \(a_{j,i}\ne a'_{j,i}\)，则
  \[
  C^{(L+1)}_{a,t}\cap C^{(L+1)}_{a',t}=\varnothing.
  \]
\end{enumerate}
因此，递归地址化中的派生概念先由有限旧层原子构造，再由包含或相交关系进行谓词化；任何 incidence 断言均可回溯到同一张有限 Boolean 账本。
\end{proposition}

% --- Distillation writeback ---
\begin{theorem}[地址歧义的条件熵降与无界寄存器泄漏]\label{distill:bourgain-entropy_increment_closure-address-ambiguity-entropy-leak}
固定状态 \(p\) 与已对象化地址 \(a\in A(p)\)。设
\[
E\subseteq \Gamma_p(a)
\]
为一族有限候选局部截面，\(\mathfrak m\) 为 \(E\) 上满支撑概率测度，\(\rho:E\to Y\) 为当前协议已经登记的有限可见读出。用自然对数定义 Shannon 熵，并记
\[
\mathsf H_\rho(E):=H_{\mathfrak m}(E\mid \rho)
=
\sum_{y\in\rho(E)}\mathfrak m(\rho^{-1}(y))H(\mathfrak m_y),
\]
其中 \(\mathfrak m_y\) 是 \(\rho^{-1}(y)\) 上的条件测度。称有限映射 \(\chi:E\to Z\) 为相对于 \(\rho\) 的授权细化判别器，若其图可由某个共同细化状态 \(q\le p\) 登记，且不改变地址对象 \(a\)。令
\[
\rho^\chi:E\longrightarrow Y\times Z,
\qquad
\rho^\chi(e):=(\rho(e),\chi(e)).
\]
则：
\begin{enumerate}
  \item 条件熵满足精确降公式
  \[
  \mathsf H_{\rho^\chi}(E)
  =\mathsf H_\rho(E)-H_{\mathfrak m}(\chi\mid \rho).
  \]
  特别地，\(\mathsf H_{\rho^\chi}(E)<\mathsf H_\rho(E)\) 当且仅当存在正 \(\mathfrak m\)-质量的 \(\rho\)-纤维使 \(\chi\) 在其上非常值。
  \item 若 \(\rho^\chi\) 在 \(\mathfrak m\) 的支撑上为终端读出，即每个正质量 \(\rho^\chi\)-纤维均为单点，则地址 \(a\) 在对应细化状态 \(q\) 上由 \(E\) 给出单值且非 \(\mathsf{NULL}\) 的局部读出证书。
  \item 设 \((E_N,\mathfrak m_N,\rho_N)_{N\ge0}\) 是同一地址 \(a\) 的一列有限歧义审计，并假设每个授权细化判别器都在每条正 \(\mathfrak m_N\)-质量的 \(\rho_N\)-纤维上为常值。置
  \[
  M_N:=\max_{y\in\rho_N(E_N)}\bigl|E_N\cap\rho_N^{-1}(y)\bigr|.
  \]
  若 \(\sup_N M_N=\infty\)，则满足如下量化泄漏性质：对每个 \(B\in\NN\)，存在 \(N\) 与一条正质量纤维 \(F\subseteq E_N\)，使得任意 \(|R|\le B\) 的有限寄存器读出 \(r:F\to R\) 都不能把 \(F\) 终端分离。
\end{enumerate}
\end{theorem}
\begin{proof}
因为 \(\chi\) 是 \(E\) 的函数，条件熵链式法则给出
\[
H_{\mathfrak m}(E\mid \rho)
=H_{\mathfrak m}(\chi,E\mid \rho)
=H_{\mathfrak m}(\chi\mid \rho)+H_{\mathfrak m}(E\mid \rho,\chi).
\]
而 \(H_{\mathfrak m}(E\mid \rho,\chi)=\mathsf H_{\rho^\chi}(E)\)，于是得到第一条公式。有限概率空间上，条件熵 \(H_{\mathfrak m}(\chi\mid \rho)\) 为零当且仅当在每条正质量 \(\rho\)-纤维上 \(\chi\) 的条件分布为 Dirac 测度，也即 \(\chi\) 为常值；严格性判据随之成立。

若 \(\rho^\chi\) 的每条正质量纤维均为单点，则每个可见读出值至多对应一个候选局部截面。由于 \(E\subseteq\Gamma_p(a)\)，且 \(q\le p\) 只是在同一地址对象上登记额外判别信息，已存在的局部截面证书不因细化而被撤销。由推论 \ref{cor:observer-indexed-nonnull-criterion}，地址已对象化且存在唯一可核验局部截面时给出非 \(\mathsf{NULL}\) 读出，第二条成立。

对第三条，任取 \(B\in\NN\)。由 \(\sup_N M_N=\infty\)，可取 \(N\) 与 \(y\in\rho_N(E_N)\) 使
\[
F:=E_N\cap\rho_N^{-1}(y),
\qquad |F|>B.
\]
若某个 \(|R|\le B\) 的寄存器读出 \(r:F\to R\) 终端分离 \(F\)，则 \(r\) 必须在 \(F\) 上单射；否则两个不同候选截面具有相同寄存器值而仍未被区分。单射性给出 \(|R|\ge |F|>B\)，与 \(|R|\le B\) 矛盾。假设中授权细化判别器均在正质量歧义纤维上为常值，故这类纤维不能通过内部细化产生第一条中的严格熵降。量化泄漏性质得证。
\end{proof}

% --- End distillation writeback ---

\begin{proof}
第一条只是把条件
\[
\mathbf b^{(L)}_{t+j}(x)=a_j
\qquad(0\le j<\ell)
\]
按坐标展开。对固定 \((j,i)\)，等式的第 \(i\) 个坐标为 \(1\) 时要求 \(x\in T^{-(t+j)}C_i^{(L)}\)，为 \(0\) 时要求 \(x\notin T^{-(t+j)}C_i^{(L)}\)。对全部 \((j,i)\) 取有限交即得构造式。

若 \(b\) 延长 \(a\)，则 \(C^{(L+1)}_{b,t}\) 的定义包含 \(C^{(L+1)}_{a,t}\) 的全部原子条件，并额外加入更深时刻的条件，故有包含关系。若 \(a\) 与 \(a'\) 在某个坐标 \((j,i)\) 不同，则一个柱事件要求 \(x\in T^{-(t+j)}C_i^{(L)}\)，另一个要求 \(x\notin T^{-(t+j)}C_i^{(L)}\)，二者不相交。证毕。
\end{proof}

\begin{proof}
由命题 \ref{prop:observer-indexed-common-refinement-decidability} 与命题 \ref{prop:observer-indexed-delayed-classification-not-retrocausal} 联立即得。
\end{proof}

\paragraph{类型化读出偏函数与对象层值}
前述 forcing 框架还可进一步压缩为一个状态化的类型读出偏函数。
它的作用是把“地址是否已对象化”“局部截面是否存在”与“值类是否已在当前类型制度中唯一确定”三件事统一到同一个对象层接口上。

\begin{definition}[类型化读出偏函数]\label{def:observer-indexed-typed-readout}
对任意信息状态 \(p\) 与类型化地址 \(a\)，定义偏函数 \(\mathrm{Read}_p(a)\) 如下：
\begin{enumerate}
  \item 若 \(a\notin A(p)\)，则 \(\mathrm{Read}_p(a)\) 在对象层不定义；
  \item 若 \(a\in A(p)\) 但 \(\Gamma_p(a)=\varnothing\)，则记
  \[
  \mathrm{Read}_p(a)=\mathsf{NULL};
  \]
  \item 若 \(a\in A(p)\) 且 \(\Gamma_p(a)\) 在当前类型制度下确定唯一可核验值类，则记该值类为 \(\mathrm{Read}_p(a)\)。
\end{enumerate}
当可见域进一步固定为 unitary slice 时，记相应特化为
\[
\mathrm{Read}_{\mathrm{US},p}(a).
\]
其跨状态的全局编译将在后续逆极限口径下组织为偏函数 \(\mathrm{Read}_{\mathrm{US}}\)，见定理 \ref{thm:recursive-addressing-readout-separatedness-null}。
\end{definition}

\begin{proposition}[\texorpdfstring{$\mathrm{Read}_{\mathrm{US}}$}{Read\_US} 的 forcing 判据]\label{prop:observer-indexed-readout-forcing-criterion}
\leanverified{paper\_recursive\_addressing\_observer\_indexed\_readout\_forcing\_criterion}
在当前类型纪律下，对任意状态 \(p\) 与类型化地址 \(a\)，有
\[
\mathrm{Read}_{\mathrm{US},p}(a)\neq \mathsf{NULL}
\]
当且仅当当前状态上存在协议允许、类型单值且可核验的局部截面。
并且若 \(q\le p\) 且 \(\mathrm{Read}_{\mathrm{US},p}(a)\neq \mathsf{NULL}\)，则
\[
\mathrm{Read}_{\mathrm{US},q}(a)\neq \mathsf{NULL},
\]
且其对象层值类与 \(\mathrm{Read}_{\mathrm{US},p}(a)\) 一致。
\end{proposition}

\begin{proof}
前半句由定义 \ref{def:observer-indexed-typed-readout} 直接给出。
对后半句，由命题 \ref{prop:observer-indexed-forcing-monotonicity}，一旦某个读出命题在较粗状态上已可断言，则在任何细化状态上仍可断言。
而定义 \ref{def:observer-indexed-typed-readout} 的第三条已把合法局部截面压缩为唯一可核验值类，因此细化只会增加证书支撑，不会改变对象层值类本身。
故结论成立。
\end{proof}

\begin{corollary}[显式提升原则的 forcing 解释]\label{cor:observer-indexed-explicit-lifting}
\leanverified{paper\_recursive\_addressing\_observer\_indexed\_explicit\_lifting}
若某个复合表达式同时调用来自不同精度索引、不同可见域或不同证书层级的读出，而系统并未提供一个共同细化状态 \(q\le p\)，使这些读出能在同一类型化域中同时解释，则该表达式在状态 \(p\) 上不定义。
\end{corollary}

\begin{proof}
复合表达式的对象层解释要求其各子项在同一可比较域中具有共同语义。
若不存在共同细化把这些子项提升到同一类型层，则不存在统一的对象层评价函数，因而该表达式不能在 \(p\) 上被强制成立。
这正是“混合精度的操作在类型层不定义”的 forcing 口径。
\end{proof}

\paragraph{值保持重写的对象层语义}
一旦地址对象已经形成、局部截面已经存在并且对象层值类已被确定，后续的稳定算术便不再处理“是否存在对象”，而只处理“如何在同一对象层值类内选取规范表示”。
相应地，值保持重写必须在对象层而非实现层被表述。

\begin{definition}[值保持重写的对象层语义]\label{def:observer-indexed-value-preserving-rewrite}
设 \(t\Rightarrow t'\) 是后续稳定算术中的一条重写规则。
称该规则在对象层上\emph{值保持}，若对任意信息状态 \(p\) 与任意使两侧都可解释的地址--值环境 \(\eta\)，均有
\[
p\Vdash \operatorname{Val}_{p,\eta}(t)=\operatorname{Val}_{p,\eta}(t').
\]
其中 \(\operatorname{Val}_{p,\eta}(t)\) 表示项 \(t\) 在状态 \(p\) 与环境 \(\eta\) 下的对象层读出类。
\end{definition}

\begin{proposition}[值保持重写不创造对象层值]\label{prop:observer-indexed-value-preserving-no-creation}
若一条重写规则在定义 \ref{def:observer-indexed-value-preserving-rewrite} 的意义下值保持，则它只能在已有对象层读出上改变表示，而不能从一个未定义项中创造出新的对象层值。
\leanverified{paper\_recursive\_addressing\_observer\_indexed\_value\_preserving\_no\_creation}
\end{proposition}

\begin{proof}
若 \(t\) 在状态 \(p\) 上未定义，则 \(\operatorname{Val}_{p,\eta}(t)\) 不存在。
此时若 \(t'\) 在 \(p\) 上已定义，则
\[
p\Vdash \operatorname{Val}_{p,\eta}(t)=\operatorname{Val}_{p,\eta}(t')
\]
无从成立，因此该规则不可能是对象层值保持的。
反之，若两侧都已定义，则值保持只说明二者属于同一对象层值类，因而改变的是表示、规范形或证书接口，而不是值本身。
\end{proof}

\begin{corollary}[规范化不创造事实]\label{cor:observer-indexed-normalization-no-fact-creation}
后续的规范化、重写与横截选择，其对象层作用只是把一个已定义值的不同表示压到统一规范形；它们不能替代地址对象化、证书建立或前缀细化，更不能越过 \(\mathsf{NULL}\) 直接生成数值事实。
\leanverified{paper_recursive_addressing_observer_indexed_normalization_no_fact_creation_seeds}
\end{corollary}

\begin{proof}
这是命题 \ref{prop:observer-indexed-value-preserving-no-creation} 的直接推论。
\end{proof}

\begin{proposition}[有限分辨率算术的对象层资格]\label{prop:observer-indexed-finite-resolution-object-eligibility}
\leanverified{paper\_recursive\_addressing\_observer\_indexed\_finite\_resolution\_object\_eligibility}
在分辨率 \(m\) 上，稳定算术只对已经通过对象化与非 \(\mathsf{NULL}\) 判据的读出类生效。
更具体地说，若
\[
\omega\equiv_m\omega'
\Longleftrightarrow
\Fold_m(\omega)=\Fold_m(\omega')
\]
给出实现层的折叠同余，则对象层算术并不直接作用于原始微态 \(\omega\)，而作用于由当前状态所支撑的读出类，即实现层等价类经地址化与类型过滤之后的对象。
\end{proposition}

\begin{proof}
折叠同余只说明不同微态在当前分辨率下不可区分；它本身既不构成对象层地址，也不自动给出可核验读出。
只有当相应等价类在当前信息状态上被地址化、通过协议并获得局部截面时，它才上升为对象层对象。
故稳定算术处理的不是原始微态，而是 forcing 已经承认的读出类。
\end{proof}

\begin{corollary}[第 \texorpdfstring{\ref{sec:emergent-arithmetic}}{7} 节的读法]\label{cor:observer-indexed-reading-of-emergent-arithmetic}
\leanverified{paper\_recursive\_addressing\_observer\_indexed\_reading\_of\_emergent\_arithmetic}
第 \ref{sec:emergent-arithmetic} 节中的值保持重写、规范横截面与分辨率模结构，应理解为对那些已经在本节意义下完成对象化、并已获得非 \(\mathsf{NULL}\) 读出的地址--值对，构造其规范表示与值保持演算。
更具体地说：
\begin{enumerate}
  \item 第 \ref{sec:emergent-arithmetic} 节并不是在无地址对象上直接施行算术；
  \item 其中的“值保持”并不是实现层字符串恒等，而是对象层读出类恒等；
  \item 其中的“模结构”也不是任意商，而是由分辨率折叠与可核验读出共同决定的对象层商结构。
\end{enumerate}
\leanverified{paper\_recursive\_addressing\_observer\_indexed\_reading\_of\_emergent\_arithmetic}
\end{corollary}

% --- Distillation writeback ---
\begin{proposition}[对象层等式的有限 forcing 账本正规形]\label{distill:euclid_elements-finite_construction_certificate_families-forcing-equality-ledger-normal-form}
\textbf{日期时间：2026-04-23T11:58:48+08:00。依赖状态：rigidity.}
固定信息状态 \(p\)。考虑由下列 primitive 语法生成的对象层等式证明：
\begin{enumerate}
  \item 已在某个共同细化状态 \(q\le p\) 上非 \(\mathsf{NULL}\) 的读出 \(\mathrm{Read}_{\mathrm{US},q}(a)\)；
  \item 定义 \ref{def:observer-indexed-value-preserving-rewrite} 意义下已登记的值保持重写；
  \item common-notion 规则：自反、对称、传递以及对已定义对象层值的替换。
\end{enumerate}
称一份有限 equality ledger 为一棵有限有向无环证明树，其叶节点只允许是同一非 \(\mathsf{NULL}\) 读出的对象层同一性或已登记值保持重写，其内节点只允许使用上述 common-notion 规则。则在该语法内：
\begin{enumerate}
  \item 若参与复合表达式的各读出没有共同细化状态使其落在同一类型化域中，则该表达式在状态 \(p\) 上不定义；
  \item 任一有限 equality ledger 的根结论
  \[
  t=t'
  \]
  都满足
  \[
  p\Vdash \operatorname{Val}_{p,\eta}(t)=\operatorname{Val}_{p,\eta}(t')
  \]
  对所有使两侧已定义的地址--值环境 \(\eta\) 成立；
  \item 反过来，凡是在上述 primitive 语法中推出的对象层等式，都可整理为一份有限 equality ledger。因此，在该语法内，等式不能由图示相似性或规范化步骤凭空产生；它必须回指到某个非 \(\mathsf{NULL}\) 读出、某条值保持重写或一条 common-notion 规则。
\end{enumerate}
\end{proposition}

% --- Distillation writeback ---
\begin{proposition}[随机分支读出的确定性冻结与提升门控]\label{prop:distill-bourgain-forcing-random-certificate-extraction}
\textbf{依赖状态：conditional audit / conditional rigidity.}
固定信息状态 \(p\)，有限指标集 \(I\)，以及对每个 \(i\in I\) 的细化状态 \(q_i\le p\)。令 \(\mathcal E\) 表示一个对象层等式断言 \(t=t'\)。假设给定有限审计数据
\[
(\lambda_i,L_i,\sigma_i,\pi_i)_{i\in I},
\]
其中 \(\lambda_i\in\QQ_{\ge0}\)、\(\sum_i\lambda_i=1\)，\(L_i\) 是一份候选有限 equality ledger 或空符号，\(\sigma_i,\pi_i\in\{0,1\}\)。这些数据的语义规定如下：若 \(\sigma_i=1\)，则 \(L_i\) 是在状态 \(q_i\) 上证明 \(\mathcal E\) 的有限 equality ledger，因而对所有使两侧在 \(q_i\) 上已定义的地址--值环境 \(\eta\) 有
\[
q_i\Vdash \operatorname{Val}_{q_i,\eta}(t)=\operatorname{Val}_{q_i,\eta}(t').
\]
若 \(\pi_i=1\)，则验证器还显式登记了提升规则
\[
(q_i,L_i,\mathcal E)\leadsto_p \mathcal E,
\]
其含义正是
\[
p\Vdash \operatorname{Val}_{p,\eta}(t)=\operatorname{Val}_{p,\eta}(t')
\]
对所有使两侧在 \(p\) 上已定义的环境 \(\eta\) 成立。定义该有限随机分支族的可提升成功质量为
\[
M_p(I):=\sum_{i\in I}\lambda_i\sigma_i\pi_i.
\]
则：
\begin{enumerate}
  \item 若 \(M_p(I)>0\)，则存在 \(i_\ast\in I\) 使 \(\lambda_{i_\ast}>0\)、\(\sigma_{i_\ast}=1\)、\(\pi_{i_\ast}=1\)。冻结三元组 \((q_{i_\ast},L_{i_\ast},\mathcal E)\) 后得到状态 \(p\) 上的一份确定性对象层等式证书；
  \item 若 \(\sum_i\lambda_i\sigma_i>0\)，则至少存在某个细化状态 \(q_i\) 上的确定性等式 ledger；但除非相应 \(π_i=1\)，该 ledger 只停留在 \(q_i\) 层，不能自动升级为 \(p\) 上的 forcing 结论；
  \item 若 \(M_p(I)=0\)，则只能推出当前抽样分支族中没有同时满足成功与提升登记的分支。它不推出 \(p\) 上 \(\mathcal E\) 为假，也不排除存在未列入 \(I\) 的其它确定性证明账本。
\end{enumerate}
\end{proposition}
\begin{proof}
第一条由有限和的正性直接给出：若每个正权分支都满足 \(\sigma_i\pi_i=0\)，则 \(M_p(I)=0\)，与假设矛盾。因此可取 \(i_\ast\) 使 \(\lambda_{i_\ast}>0\) 且 \(\sigma_{i_\ast}=\pi_{i_\ast}=1\)。由 \(\sigma_{i_\ast}=1\)，\(L_{i_\ast}\) 是 \(q_{i_\ast}\) 上的有限 equality ledger；再由 \(\pi_{i_\ast}=1\) 的定义，登记的提升规则把该 ledger 编译为状态 \(p\) 上的对象层等式证书。

第二条同理，\(\sum_i\lambda_i\sigma_i>0\) 给出一个正权成功分支，但该结论只使用 \(\sigma_i\) 的语义，因此只得到 \(q_i\) 层 forcing。forcing 的通常细化单调性并不提供从 \(q_i\le p\) 回到 \(p\) 的反向提升；反向提升恰由 \(\pi_i=1\) 单独登记。

第三条只是对有限审计域的否定陈述。\(M_p(I)=0\) 说明在已列出的有限分支中没有正权分支同时通过成功位与提升位；该命题没有量化所有可能 ledger，故不能推出对象层断言在 \(p\) 上不可证或为假。证毕。
\end{proof}

% --- End distillation writeback ---

\begin{proof}
第一条正是推论 \ref{cor:observer-indexed-explicit-lifting}：复合表达式要求所有子项在同一可比较域中解释；若不存在共同细化状态，则没有统一对象层评价函数，故表达式不定义。

第二条对 ledger 的高度作归纳。高度为零时，叶节点分两类。若叶节点是同一非 \(\mathsf{NULL}\) 读出的对象层同一性，则由定义 \ref{def:observer-indexed-typed-readout} 的第三条，读出类已经唯一确定，故等式成立。若叶节点是已登记的值保持重写，则结论正是定义 \ref{def:observer-indexed-value-preserving-rewrite}。归纳步中，自反、对称与传递保持对象层等式；替换规则只在已定义对象层值上使用，因而由等式的相容性保持强制成立。故 ledger 根结论被 \(p\) 强制。

第三条由语法本身的有限生成性给出。任一推导只有有限多步；把其原始读出同一性与值保持重写作为叶节点，把每次 common-notion 应用作为内节点，即得到有限有向无环证明树。命题 \ref{prop:observer-indexed-value-preserving-no-creation} 进一步保证值保持重写不能从未定义项创造对象层值；因此任何合法等式推导都不可能越过 \(\mathsf{NULL}\) 或未构造地址而生成事实。证毕。
\end{proof}

% --- End distillation writeback ---

\begin{proof}
由推论 \ref{cor:observer-indexed-nonnull-criterion}，对象层算术的合法输入首先必须是非 \(\mathsf{NULL}\) 的地址--值对；
由命题 \ref{prop:observer-indexed-value-preserving-no-creation}，值保持重写只能在已定义对象上改变表示；
再由命题 \ref{prop:observer-indexed-finite-resolution-object-eligibility}，有限分辨率下真正进入算术层的对象，是折叠同余经对象化与类型过滤之后的读出类。
三者合并即得结论。
\end{proof}

\paragraph{本节收束}
至此，第 \ref{sec:recursive-addressing} 节可压缩为一条统一逻辑链：
先由地址生成把对象带入可指称域；
再由前缀站点、局部截面与证书对象决定何者可读；
再由 forcing 语义区分成立、失败与未决定；
再由多轴前缀与显式提升把未决定命题推进到可断言；
最后，只有在这一切都完成之后，稳定算术才获得其合法适用域。
因此，第 \ref{sec:emergent-arithmetic} 节并不是从零开始引入一个外加层，而是本节对象层语义的代数化延续：
本节回答“何时可以谈值”，而第 \ref{sec:emergent-arithmetic} 节回答“在已经可以谈值之后，如何进行值保持的规范化与算术化”。

